\section{Mathematical formulation}
\label{sec:governing}

The formulation is built in the order in which the constitutive statement
generates the equations of motion: the geometry and the fields are fixed first
(\cref{sec:governing_variational}), then the kinematics and the gradient energy
(\cref{sec:governing_kinematics}), then the enthalpy from which every
constitutive relation is differentiated (\cref{sec:governing_constitutive}),
then the balance laws obtained by independent variations
(\cref{sec:governing_balance}), then the electric model with the two ratios
that justify it (\cref{sec:governing_electric}), the boundary conditions
(\cref{sec:governing_bc}), the power flux that defines a backward mode
(\cref{sec:governing_energy}), the dimensionless groups
(\cref{sec:governing_nondim}), the limits of the model
(\cref{sec:governing_limits}) and the condition under which the enthalpy stays
positive definite (\cref{sec:governing_stability}).

\subsection{Geometry, fields and notation}
\label{sec:governing_variational}

Consider a transversely isotropic ($6mm$) n-type ZnO plate of thickness $d$
whose faces are the planes $z=\pm d/2$, occupying the half-space $x<0$ and
terminating at a traction-free edge at $x=0$.  The plate is unbounded in $y$, so
a single tangential wavenumber $k_y$ is conserved, and every field has the form
$f(x,z)\,e^{ik_yy-i\omega t}$ with three displacement components retained.  The
formulation is therefore two-and-a-half-dimensional (2.5D): a discrete
cross-section in $(x,z)$ carrying full three-component kinematics, a single
conserved tangential wavenumber, and no discretisation of $y$.  The independent fields are
the displacement $u_i(\vect{x},t)$, the electrostatic potential
$\varphi(\vect{x},t)$ and the electron-number perturbation
$\Delta n(\vect{x},t)$ measured from the uniform doping $n_0$.  The polar $c$
axis of the crystal is $x_3\equiv z$; Latin indices run over $1,2,3$, a comma
denotes a spatial derivative and a superimposed dot a time derivative.

Three conventions are fixed once and used everywhere.  The electric field is
$E_i=-\partial_i\varphi$.  The electron charge is $-q$ with $q>0$.  The
quantity $\widetilde\mu$ is the \emph{electrochemical potential of an electron
per carrier}, measured in joules, so that every term appearing in it is an
energy and no term is a dimensionless ratio.

\subsection{Kinematics and the strain gradient}
\label{sec:governing_kinematics}

The small-strain tensor and its gradient are
\begin{equation}
\label{eq:strain}
\varepsilon_{ij}=\tfrac12\bigl(u_{i,j}+u_{j,i}\bigr),
\qquad
\eta_{ijk}\equiv\varepsilon_{ij,k}=\tfrac12\bigl(u_{i,jk}+u_{j,ik}\bigr),
\end{equation}
so that $\eta_{ijk}$ is symmetric in its first two indices,
$\eta_{ijk}=\eta_{jik}$, and carries the dimension of an inverse length.
All three displacement components are retained; the strain has six independent
components and the strain gradient eighteen.

The higher-order elastic energy is taken in the single-length
(isotropic-gradient) form
\begin{equation}
\label{eq:gradient_energy}
w_g=\tfrac12\,A_{ijklmn}\,\eta_{ijk}\,\eta_{lmn}
=\tfrac12\,\ell^{2}\,c_{ijlm}\,\eta_{ijk}\,\eta_{lmk},
\qquad
A_{ijklmn}=\ell^{2}\,c_{ijlm}\,\delta_{kn},
\end{equation}
where in \eqref{eq:gradient_energy} the \emph{third and sixth slots} of
$A_{ijklmn}$ are the gradient directions: the third slot is contracted with the
gradient index of the first strain gradient and the sixth with that of the
second.  This is the only contraction for which the energy depends on the
gradient of the strain through $c_{ijlm}\eta_{ijk}\eta_{lmk}$; contracting any
other pair of slots would break the major symmetry
$A_{ijklmn}=A_{lmnijk}$ that a second derivative of an energy must possess.  The two candidate contractions are also distinguishable by their spectra: the
fundamental shear-horizontal branch of \eqref{eq:gradient_energy} has the closed
form recorded in \cref{sec:numerics_verification} and is reproduced to
$2\times10^{-13}$, while the alternative contraction misses that branch by a
residual of $3\times10^{6}$.

The units follow from \eqref{eq:gradient_energy} and are collected in
\cref{tab:units}.  Because $\eta$ carries an inverse length, the gradient
modulus $A$ is \emph{not} a stiffness: it has the units of a force
($\mathrm{N}=\mathrm{J/m}$), one factor of length times the elastic stiffness.
Correspondingly the hyperstress conjugate to $\eta$ has the units of a force per
unit length ($\mathrm{N/m}$), and its double divergence has the units of a force
density, as the balance law requires.  A strain-gradient energy of the form \eqref{eq:gradient_energy} therefore
assigns a force, and not a stiffness, to $A$, and a force per unit length, and
not a stress, to the hyperstress.

\begin{table}[t]
\centering
\caption{Dimensions of the quantities of the formulation in SI base units.
The strain-gradient sector is the one that is easy to mis-state: $A$ is a force,
not a stiffness, and the hyperstress is a force per unit length, not a stress.}
\label{tab:units}
\small
\begin{tabular}{@{}llll@{}}
\toprule
Quantity & Symbol & SI unit & Base units \\
\midrule
displacement & $u_i$ & m & m \\
strain & $\varepsilon_{ij}$ & -- & 1 \\
strain gradient & $\eta_{ijk}$ & m$^{-1}$ & m$^{-1}$ \\
elastic stiffness & $c_{ijkl}$ & Pa & kg\,m$^{-1}$s$^{-2}$ \\
strain-gradient modulus & $A_{ijklmn}$ & N & kg\,m\,s$^{-2}$ \\
Cauchy stress & $\sigma_{ij}$ & Pa & kg\,m$^{-1}$s$^{-2}$ \\
hyperstress & $\tau_{ijk}$ & N\,m$^{-1}$ & kg\,s$^{-2}$ \\
double traction & $t_{ij}=n_k\tau_{ijk}$ & N\,m$^{-1}$ & kg\,s$^{-2}$ \\
piezoelectric constant & $e_{kij}$ & C\,m$^{-2}$ & A\,s\,m$^{-2}$ \\
flexoelectric constant & $f_{ijkl}$ & C\,m$^{-1}$ & A\,s\,m$^{-1}$ \\
permittivity & $\kappa_{ij}$ & F\,m$^{-1}$ & A$^2$s$^4$kg$^{-1}$m$^{-3}$ \\
carrier density & $\Delta n$ & m$^{-3}$ & m$^{-3}$ \\
carrier compressibility & $a=k_BT/n_0$ & J\,m$^3$ & kg\,m$^4$s$^{-2}$ \\
coupling constant & $H_{jkl}$ & J\,m & kg\,m$^3$s$^{-2}$ \\
electrochemical potential & $\widetilde\mu$ & J & kg\,m$^2$s$^{-2}$ \\
\bottomrule
\end{tabular}
\end{table}

\subsection{Enthalpy and constitutive relations}
\label{sec:governing_constitutive}

The internal-energy density of the model is
\begin{equation}
\label{eq:enthalpy}
\begin{aligned}
w={}&\tfrac12\,\varepsilon_{ij}c_{ijkl}\varepsilon_{kl}
+\tfrac12\,A_{ijklmn}\eta_{ijk}\eta_{lmn}
-\tfrac12\,E_i\kappa_{ij}E_j
-e_{kij}E_k\varepsilon_{ij}
-f_{ijkl}E_i\eta_{jkl}\\
&+\tfrac12\,a\,(\Delta n)^2
+H_{jkl}\,\Delta n\,\eta_{jkl}
-q\,\Delta n\,\varphi ,
\end{aligned}
\end{equation}
with $a=k_BT/n_0$.  The terms are, in order, the elastic, strain-gradient,
dielectric, piezoelectric and flexoelectric energies of the classical
flexoelectric dielectric \cite{shen2010theory}, the carrier compressibility, the
added carrier--strain-gradient coupling and the electrostatic energy of the
mobile charge.  The seventh term is the addition of this article.

Differentiating \eqref{eq:enthalpy} gives the constitutive relations
\begin{align}
\sigma_{ij}&=\frac{\partial w}{\partial\varepsilon_{ij}}
=c_{ijkl}\varepsilon_{kl}-e_{kij}E_k,
\label{eq:sigma}\\
\tau_{ijk}&=\frac{\partial w}{\partial\eta_{ijk}}
=A_{ijklmn}\eta_{lmn}+H_{ijk}\Delta n-f_{lijk}E_l,
\label{eq:tau}\\
D_i&=-\frac{\partial w}{\partial E_i}
=\kappa_{ij}E_j+e_{ijk}\varepsilon_{jk}+f_{ijkl}\eta_{jkl},
\label{eq:D}\\
\widetilde\mu&=\frac{\partial w}{\partial\Delta n}
=a\,\Delta n-q\varphi+H_{jkl}\eta_{jkl}.
\label{eq:mu}
\end{align}
Equation \eqref{eq:mu} is the electrochemical potential of an electron in
joules: the compressibility term $k_BT\,\Delta n/n_0$ is the chemical part, the
$-q\varphi$ term is the electrostatic energy of a particle of charge $-q$, and
the $H$ term is the energy the added coupling assigns to a carrier sitting in a
strain gradient.  Every term of \eqref{eq:mu} is an energy, which is the convention used throughout; a dimensionless form obtained by dividing
\eqref{eq:mu} by $k_BT$ is used internally and is recorded in
\cref{sec:governing_nondim}.

The three non-classical entries of \eqref{eq:tau}--\eqref{eq:mu} are reciprocal.
The pair $(\Delta n,\eta)$ contributes $H_{jkl}\Delta n$ to the hyperstress and
$H_{jkl}\eta_{jkl}$ to the electrochemical potential; the pair $(E,\eta)$
contributes $-f_{lijk}E_l$ to the hyperstress and $+f_{ijkl}\eta_{jkl}$ to the
electric displacement.  In both cases the same tensor appears in the two
conjugate relations, which is the content of the Maxwell identities
$\partial^2w/\partial\eta\partial(\Delta n)=\partial^2w/\partial(\Delta
n)\partial\eta$ and $\partial^2w/\partial\eta\partial E=\partial^2w/\partial
E\partial\eta$.  The reciprocity of the discretised operator is verified
numerically at $\pm k_y$ to $10^{-13}$ (\cref{sec:numerics_verification}).

The flexoelectric sector is the symmetric one generated by the single term
$-f_{ijkl}E_i\eta_{jkl}$ of \eqref{eq:enthalpy}: the direct effect appears in
$D$ and the converse effect in $\tau$, both through the same constant.  In the
executed calculations a single component is retained, $f\equiv f_{3zz}$, so
that $D_z=e_{zjk}\varepsilon_{jk}+\kappa_{zz}E_z+f\,\eta_{zzz}$ and
$\tau_{zzz}=A_{zzz\,lmn}\eta_{lmn}-f E_z$.  Neither the surface flexoelectric
effect of \cite{shen2010theory} nor the electric-field-gradient (higher-order)
flexoelectric terms are included.  The value of $f$ used, its provenance and
its measured influence on the reported numbers are stated in
\cref{sec:results_scope}.

\subsection{Balance laws}
\label{sec:governing_balance}

Stationarity of the action with respect to $u_j$ gives the mechanical balance
\begin{equation}
\label{eq:momentum}
\rho\,\ddot u_j=\partial_i\sigma_{ij}-\partial_i\partial_k\tau_{jik},
\end{equation}
where the double divergence of the hyperstress has the sign that makes the
gradient energy energetically favourable to smooth fields: the variation of
$\tfrac12A(\cdot)^2$ produces $+\partial_i\partial_k(A\,\eta)$ on the left of
the weak form and therefore a minus sign here.  It is convenient to define the
\emph{effective stress}
\begin{equation}
\label{eq:effective}
\Sigma_{ij}=\sigma_{ij}-\partial_k\tau_{ijk},
\end{equation}
in terms of which \eqref{eq:momentum} is the familiar $\rho\ddot u_j=
\partial_i\Sigma_{ij}$, and the traction on a surface is $t_j=n_i\Sigma_{ij}$.
$\Sigma$ is not the quantity that enters the constitutive law; it is the
divergence of the hyperstress that makes \eqref{eq:momentum} look classical.

Variation with respect to $\varphi$ gives Gauss's law
\begin{equation}
\label{eq:gauss}
\partial_iD_i=-q\,\Delta n,
\end{equation}
the minus sign being the charge of the electron, and variation with respect to
$\Delta n$ gives the carrier continuity equation in electrochemical form,
\begin{equation}
\label{eq:continuity}
q\,\partial_t(\Delta n)+\partial_iJ_i=0,
\qquad
J_i=-\frac{D\,n_0}{k_BT}\,\partial_i\widetilde\mu ,
\end{equation}
where $J_i$ is the electron \emph{number} flux and $D$ the diffusivity.  The
transport coefficient of \eqref{eq:continuity} follows from the Einstein
relation: with $\mu_e$ the mobility, $D=\mu_ek_BT/q$, so that
$Dn_0/k_BT=\mu_en_0/q$ and \eqref{eq:continuity} carries exactly the
drift--diffusion content.  The associated charge-current density is
\begin{equation}
\label{eq:current}
j_i=-qJ_i=\frac{qDn_0}{k_BT}\,\partial_i\widetilde\mu
=qD\,\partial_i\Delta n-q\mu_e n_0\,\partial_i\varphi
+\frac{qDn_0}{k_BT}H_{jkl}\,\partial_i\eta_{jkl},
\end{equation}
in which the three terms are the diffusion current, the drift current and the
carrier--strain-gradient contribution, with the familiar electron signs: the
diffusion current flows up the density gradient and the drift current against
the field.  The last term is a prediction of the added coupling: a
strain-gradient field drives a current even at zero applied field, which is the
transport signature of the term \eqref{eq:added_coupling}.

\begin{table}[t]
\centering
\caption{Conjugate pairs of the enthalpy \eqref{eq:enthalpy}, the constitutive
relation generated by each and its dimension.  The last two rows are the only
non-classical entries.}
\label{tab:conjugates}
\begin{tabular}{@{}lllll@{}}
\toprule
Generalised displacement & Force & Relation & Unit of force & Origin \\
\midrule
$\varepsilon_{ij}$ & $\sigma_{ij}$ & \cref{eq:sigma} & Pa & elasticity \\
$\eta_{ijk}$ & $\tau_{ijk}$ & \cref{eq:tau} & N\,m$^{-1}$ & gradient + $H$ + flexo \\
$E_i$ & $D_i$ & \cref{eq:D} & C\,m$^{-2}$ & dielectric + piezo + flexo \\
$\Delta n$ & $\widetilde\mu$ & \cref{eq:mu} & J & compressibility + $H$ \\
\bottomrule
\end{tabular}
\end{table}

\subsection{The electric model and the quasistatic approximation}
\label{sec:governing_electric}

Equations \eqref{eq:gauss}--\eqref{eq:continuity} are electrostatic: no
magnetic induction, no displacement-current term in Amp\`ere's law and no
electromagnetic radiation mode are retained.  The approximation is justified at the operating point by evaluating the two
ratios that control it.

\emph{Retardation.}  At $f=\SI{16.68}{\giga\hertz}$ ($\omega=\SI{1.048e11}{%
\per\second}$) and with the relative permittivity $\kappa_{33}/\varepsilon_0=10$
of ZnO, the electromagnetic wavenumber in the medium is
$k_{\mathrm{EM}}=\omega\sqrt{\kappa_r}/c=\SI{1.106e3}{\per\metre}$, so that
$\lambda_{\mathrm{EM}}=2\pi/k_{\mathrm{EM}}=v_{\mathrm{EM}}/f=\SI{5.68}{%
\milli\metre}$ while $1/k_{\mathrm{EM}}=\SI{0.905}{\milli\metre}$.  The plate is
$d=\SI{200}{\nano\metre}$ thick, hence $d/\lambda_{\mathrm{EM}}=3.52\times10^{-5}$
and $k_{\mathrm{EM}}d=2.2\times10^{-4}$; the retardation correction is
controlled by $k_{\mathrm{EM}}d$ and not by $d/\lambda_{\mathrm{EM}}$, and it is
of order $(k_{\mathrm{EM}}d)^2=4.9\times10^{-8}$.  The wave itself gives the same conclusion: its phase speed is
$v_p=(\omega/|\vect{k}|)=\SI{28.0}{\kilo\metre\per\second}$, i.e.\
$v_p/c=9.4\times10^{-5}$, so the magnetic energy is smaller than the electric
energy by $(v_p/c)^2\approx8.8\times10^{-9}$.  The quasistatic reduction is therefore accurate, to the order quoted, at the
reported operating points.

\emph{Carrier relaxation.}  The opposite conclusion holds for the carriers, and
this is the substantive point.  With $\sigma_e=q\mu_en_0=\SI{5.6e-3}{\siemens
\per\metre}$ the dielectric relaxation time is
$\tau_d=\kappa_{33}/\sigma_e=\SI{1.58e-8}{\second}$, whereas the wave period is
$T=2\pi/\omega=\SI{6.0e-11}{\second}$; hence $\omega\tau_d=1.65\times10^{3}$.
The carriers cannot follow the field quasistatically at this frequency.  The
model therefore does not replace charge relaxation by an instantaneous
equilibrium: the full continuity equation \eqref{eq:continuity} is retained and
the carrier response is part of the eigenvalue problem.  What is neglected is
electromagnetic retardation, not charge dynamics.

\emph{Screening.}  The screening length is $L_D=\sqrt{\kappa_{33}k_BT/(q^2n_0)}
=\SI{2.68}{\micro\metre}$, so $d/L_D=0.075$: the plate is much thinner than the
Debye length and the carrier perturbation is not confined to a surface layer.
All three ratios are computed from the material file and stored as part of the
verification record.

\subsection{Boundary conditions}
\label{sec:governing_bc}

The boundary conditions follow from the surface terms of the variational
principle.  They are stated separately for the two faces, the reflecting edge
and the far field, because the physics of each is different.

\paragraph{Plate faces $z=\pm d/2$.}  The faces are free surfaces, so the
traction and the double traction vanish,
\begin{equation}
\label{eq:bc_faces_mech}
\sigma_{iz}=0,\qquad \tau_{ijz}=0,\qquad z=\pm d/2,
\end{equation}
the electrons cannot leave the crystal (blocking contacts, no surface
recombination),
\begin{equation}
\label{eq:bc_faces_carrier}
J_z=0,\qquad z=\pm d/2,
\end{equation}
and the electrical condition is one of
\begin{equation}
\label{eq:bc_grounded}
\varphi=0\quad\text{(grounded)},\qquad
D_z=0\quad\text{(open circuit)},\qquad z=\pm d/2 .
\end{equation}
The grounded condition models perfectly conducting electrodes on both faces
shorted to each other; it is \emph{essential} and is imposed on the basis
functions.  The open-circuit condition models a bare surface with no free
charge; it is \emph{natural} and is imposed weakly.  In the open-circuit case
the potential is determined by \eqref{eq:gauss} only up to an additive constant,
which is a genuine null vector of the electrostatic operator and is removed by a
zero-thickness-mean gauge.  The gauge changes no field, no charge density and no
power flux, but it also means that the open-circuit results are a solvable
version of the same physics rather than an independent validation of it.

\paragraph{Reflecting edge $x=0$.}  The edge is a free surface in vacuum:
\begin{equation}
\label{eq:bc_edge}
\sigma_{ix}=0,\qquad \tau_{ijx}=0,\qquad D_x=0,\qquad J_x=0,\qquad x=0 .
\end{equation}

\paragraph{Far field $x\to-\infty$.}  The solution is an incoming guided mode
plus outgoing branches; no wave arrives from $x=-\infty$ other than the
prescribed incident one.

\paragraph{The corner.}  Where the edge meets a face the traction and double
traction conditions are not independent: the classical Mindlin decomposition
produces a line force along the corner whose density is
$n_z\tau_{ijz}|_{x=0}$ evaluated on the limit of the face conditions.  That line
force is a weak singularity of the discretisation and is \emph{not} resolved
pointwise here; the weak form is imposed exactly at the corners, the line force is non-zero
there ($>10^{-4}$), and \cref{sec:disc_limitations} shows that the reported
partitions do not rest on the energy it contains.  The geometry solved is
therefore a plate with a straight edge and a corner field that is not
resolved.

\subsection{Power flux and the energy identity}
\label{sec:governing_energy}

The power flowing in the $x$ direction, per unit width and integrated over the
thickness, is
\begin{equation}
\label{eq:Px_int}
\mathcal{P}_x=\tfrac12\Real\int_{-d/2}^{d/2}
\Bigl[-\sigma_{xi}\dot u_i^{*}-\tau_{jkx}\dot\eta_{jk}^{*}
+\varphi\dot D_x^{*}+\widetilde\mu\,\dot\Gamma_x^{*}\Bigr]\dd z ,
\end{equation}
where $\Gamma_x$ is the $x$ component of the carrier flux; the four terms are
the elastic, hyperelastic, dielectric and chemical contributions, each with the
dimensions of a power per unit area ($\mathrm{W\,m^{-2}}$); that the four are
mutually consistent is verified by an independent dimensional check.  It is this
quantity, and not the magnitude of the wavevector, that decides whether a mode
is backward: a mode is backward when $\mathcal{P}_x<0$ while it is outgoing at
the edge.

For a scattering solution the identity
\begin{equation}
\label{eq:energy_identity}
\underbrace{P_{\text{inc}}}_{\text{incident}}
+\underbrace{\sum_{j}P_j}_{\text{self-powers}}
+\underbrace{P_{\text{cross}}}_{\text{interference}}
=\underbrace{P_{\text{edge}}-P_{\text{diss}}}_{\text{flux minus loss}}
\end{equation}
holds exactly, with $P_j$ the \emph{signed} self-power of outgoing mode $j$,
$P_{\text{cross}}$ the bilinear interference between them (non-zero whenever the
modes are not orthogonal in the power metric) and
\begin{equation}
\label{eq:dissipation}
P_{\text{diss}}=\tfrac12\alpha_{\mathrm{d}}\iint
\Bigl[(k_x^2+k_y^2)|\widetilde\mu|^2+|\partial_z\widetilde\mu|^2\Bigr]\dd z ,
\qquad \alpha_{\mathrm{d}}=\frac{D}{d\,c_s},
\end{equation}
written in the scaled variables of \cref{sec:governing_nondim} --- the wavenumber
components and the thickness coordinate are the scaled ones and
$\widetilde\mu$ is the electrochemical potential divided by $k_BT$, which makes
$\alpha_{\mathrm{d}}$ the dimensionless ratio of the diffusion rate to the
elastic rate --- and it is the conduction loss inside the control volume.  The
symbol $\alpha$ itself is used later for the spatial attenuation
$2|\Imag k_x|$ of a mode (\cref{sec:numerics}); the group
$\alpha_{\mathrm{d}}$ is a property of the material and of the thickness, not of
a mode.  The terms of
\eqref{eq:energy_identity} are not independent energy channels: only the
propagating members of $\{P_j\}$ transport energy to infinity, the evanescent
members are reactive at the plane, and $P_{\text{cross}}$ is an interference
correction.  The residual of \eqref{eq:energy_identity} is therefore evaluated for every
reported solution and is plotted in \cref{fig:fig05}(c).

\subsection{Non-dimensionalisation}
\label{sec:governing_nondim}

The calculation uses lengths scaled by $d$, stresses by $c_{44}$, time by
$d/c_s$ with $c_s=\sqrt{c_{44}/\rho}$, carrier densities by $n_0$ and energies
per carrier by $k_BT$.  Two derived scales appear repeatedly,
\begin{equation}
\label{eq:a_scale}
a=\frac{k_BT}{n_0},\qquad L_D=\sqrt{\frac{\kappa_{33}k_BT}{q^2n_0}},
\end{equation}
and the model is governed by the five dimensionless numbers of
\cref{tab:groups}.  Each is an exact ratio of two physical lengths, times or
energies, and each is reproduced numerically from the material file in
\cref{sec:numerics_verification}.

\begin{table}[t]
\centering
\caption{Dimensionless groups, their definitions and their values at the
operating point ($d=\SI{200}{\nano\metre}$,
$n_0=\SI{2e18}{\per\cubic\metre}$,
$\mu_e=\SI{0.0175}{\square\metre\per\volt\per\second}$, $\ell=\SI{50}{\nano\metre}$,
$T=\SI{301.7}{\kelvin}$).  The temperature is not chosen independently: it is
the value implied by the Einstein relation for the diffusivity--mobility pair
used, $T=qD/(\mu_ek_B)$.}
\label{tab:groups}
\begin{tabular}{@{}llll@{}}
\toprule
Group & Definition & Value & Meaning \\
\midrule
$Q$ & $d/L_D$ & $7.46\times10^{-2}$ & thickness over Debye length \\
$\alpha_{\mathrm{d}}$ & $D/(d\,c_s)$ & $8.28\times10^{-1}$ & diffusion rate over elastic rate \\
$\lambda$ & $\ell/d$ & $2.50\times10^{-1}$ & gradient length over thickness \\
$f$ & $f_{3zz}/(d\sqrt{c_{44}\kappa_{33}})$ & $2.56\times10^{-3}$ & flexoelectric coupling \\
$h$ & $H/(d\sqrt{a\,c_{44}})$ & $\eta\times2.12\times10^{-1}$ & carrier--gradient coupling \\
\bottomrule
\end{tabular}
\end{table}

The coupling constant $h$ of \cref{tab:groups} is dimensionless for the
following reason.  The compressibility $a=k_BT/n_0$ has units of J\,m$^3$ and
$c_{44}$ of J\,m$^{-3}$, so $\sqrt{a\,c_{44}}$ has units of J and
$d\sqrt{a\,c_{44}}$ has units of J\,m --- exactly the units of $H$ required by
\eqref{eq:enthalpy}.  The same identity fixes the reference value
$H_{\mathrm{ref}}=d\sqrt{a\,c_{44}}=\SI{1.89e-21}{\joule\metre}$ used to
non-dimensionalise the coupling in the tables.

\subsection{Limits of the formulation}
\label{sec:governing_limits}

The model of \eqref{eq:enthalpy} contains four sectors that can be switched off
independently, and the limits they generate are the ones against which the
calculation is read.

\emph{Classical piezoelectric plate.}  Setting $\ell=f=H=0$ and removing the
carrier sector ($\Delta n\equiv0$) leaves the classical piezoelectric
formulation, with a polynomial in $k_x$ of degree two rather than four.

\emph{Strain-gradient dielectric.}  With $H=0$ and no carriers the enthalpy is
the flexoelectric dielectric enthalpy of \cite{shen2010theory} completed by the
gradient energy \eqref{eq:gradient_energy}.  In this limit the fundamental
shear-horizontal branch has the closed form
$\omega=\sqrt{c_{66}K^2(1+\ell^2K^2)}$, which is used as one of the checks of
\cref{sec:numerics_verification}.

\emph{Piezoelectric semiconductor without the gradient sector.}  Setting
$\ell=f=H=0$ and retaining the carriers leaves a drift--diffusion plate whose
sagittal spectrum is examined in \cref{sec:results_mechanism}: it contains no
outgoing mode with a positive axial wavenumber and a negative axial power
apart from the specular reflection.  The backward branch is therefore a property
of the gradient sector, and the negative-reflection problem is not well posed
without it.

\emph{Uncoupled carriers.}  With $H=0$ and the strain-gradient sector retained,
the carrier dynamics still enter through Gauss's law and through the continuity
equation \eqref{eq:continuity}, so the plate remains a piezoelectric
semiconductor; only the direct carrier--gradient enthalpy of
\eqref{eq:added_coupling} is removed.  This is the limit in which the controls
of \cref{sec:results_mechanism} are computed, and it is the limit that shows the
partition to be a property of the branch rather than of the added term.

Two further reductions are used in the text rather than tabulated here: the
electrostatic reduction of \cref{sec:governing_electric}, whose two controlling
ratios are evaluated at the operating point, and the thin-plate limit
$d/L_D=0.075$, in which the carrier perturbation is not confined to a surface
layer.

\subsection{Thermodynamic admissibility}
\label{sec:governing_stability}

The enthalpy \eqref{eq:enthalpy} is a quadratic form in
$(\varepsilon,\eta,E,\Delta n)$ and is positive definite only if the carrier and
strain-gradient blocks are not over-coupled.  Eliminating $\Delta n$ blockwise
gives the Schur condition
\begin{equation}
\label{eq:schur}
s\equiv\frac{H^2\,(c^{-1})_{33}}{\ell^2\,a}<1 ,
\end{equation}
where $(c^{-1})_{33}$ is the $33$ entry of the inverse of the full elastic
matrix, i.e.\ the reduced modulus
$c^{*}=c_{33}-2c_{13}^2/(c_{11}+c_{12})=\SI{154.4}{\giga\pascal}$ and not $c_{33}$
itself.  The bound is attained at
$H_{\text{bound}}=\ell\sqrt{a\,c^{*}}=\SI{8.97e-22}{\joule\metre}$, which is the
value at which the coupled stiffness of the plate ceases to be positive for
every strain-gradient field.  Every calculation reported below is confined to \eqref{eq:schur}.  All reported solutions have
$|H|\le\SI{2.0e-23}{\joule\metre}=0.022\,H_{\text{bound}}$, i.e.\ $s\le
5\times10^{-4}$: the enthalpy is strictly convex with a margin of more than two
thousand, and no reported result sits near a constitutive instability.

The figures normalise the coupling by
$H^{*}=H_{\text{bound}}/\sqrt5=\SI{4.01e-22}{\joule\metre}$, which is the value
at which $s=0.2$; the normalisation is a plotting convention with no physical
content beyond fixing the swept interval $|H/H^{*}|\le0.05$, that is
$|H|\le2.0\times10^{-23}\,\si{\joule\metre}$, the tabulated extreme.  The
separate symbol $H_{\mathrm{ref}}=d\sqrt{a\,c_{44}}=\SI{1.89e-21}{%
\joule\metre}$ used in \cref{tab:configurations} is the thickness-based
reference of \eqref{eq:enthalpy}; it is not the plotting normalisation, and the two differ.  No strain-gradient elasticity measurement for
ZnO fixes any of $\ell$, $f$ or $H$, so the reported intervals are modelling
choices and carry no experimental uncertainty; the numerical uncertainty of a
reported $R_B$ is its refinement spread, below $10^{-6}$ in relative terms
(\cref{tab:refinement}).
